% ============================================================
%  Companion deck: 50 logical qubits & 2D quantum dots on Magne
%  VQE-UCCSD initial state + Rodeo  (EIGEN-Q application)
% ============================================================
\documentclass[aspectratio=169,11pt]{beamer}

\usetheme{default}
\usecolortheme{default}
\useinnertheme{rectangles}
\setbeamertemplate{navigation symbols}{}
\setbeamertemplate{itemize items}[circle]
\setbeamertemplate{itemize subitem}[triangle]

\usepackage[T1]{fontenc}
\usepackage[utf8]{inputenc}
\usepackage{amsmath,amssymb}
\usepackage{booktabs}
\usepackage{array}

\definecolor{qnblue}{RGB}{0,60,120}
\definecolor{qnlight}{RGB}{220,232,245}
\definecolor{qngray}{RGB}{90,90,100}
\definecolor{qnaccent}{RGB}{200,80,30}
\definecolor{qngreen}{RGB}{0,120,80}

\setbeamercolor{structure}{fg=qnblue}
\setbeamercolor{frametitle}{fg=qnblue}
\setbeamercolor{title}{fg=qnblue}
\setbeamercolor{block title}{bg=qnblue,fg=white}
\setbeamercolor{block body}{bg=qnlight,fg=black}
\setbeamercolor{block title alerted}{bg=qnaccent,fg=white}
\setbeamercolor{block body alerted}{bg=qnlight,fg=black}
\setbeamerfont{frametitle}{series=\bfseries}
\setbeamerfont{title}{series=\bfseries}
\setbeamerfont{block title}{series=\bfseries}

\setbeamertemplate{frametitle}{%
  \vspace{6pt}\hspace{-2pt}%
  {\usebeamerfont{frametitle}\usebeamercolor[fg]{frametitle}\insertframetitle}\\[-3pt]%
  {\color{qnblue}\rule{\linewidth}{0.8pt}}%
}
\setbeamertemplate{footline}{%
  \hbox{\hspace*{0.06\paperwidth}%
  \color{qngray}\scriptsize EIGEN-Q $\cdot$ 50 logical qubits \& 2D quantum dots on Magne
  \hfill \insertframenumber/\inserttotalframenumber\hspace*{0.06\paperwidth}}\vspace{4pt}}

\newcommand{\ket}[1]{|#1\rangle}

\title{\large Can 50 Logical Qubits Run Serious\\ Quantum-Dot Benchmarks on Magne?}
\subtitle{\normalsize VQE--UCCSD initial states feeding the Rodeo algorithm}
\author{\small Feasibility assessment for M.~Hjorth-Jensen \& the EIGEN-Q team}
\date{\small June 2026}

\begin{document}

\begin{frame}[plain]
  \titlepage
  \begin{center}\small\color{qngray}
  $N$ electrons in a 2D harmonic trap $+$ Coulomb $\cdot$
  HO basis of $K$ spin-orbitals $\cdot$ Jordan--Wigner
  \end{center}
\end{frame}

% ---------- BLUF ----------
\begin{frame}{Bottom line up front}
\begin{block}{The one-sentence answer}
For 2D quantum dots, \textbf{qubit count is not the bottleneck} --- 50 logical
qubits comfortably hold the largest natural basis ($K=42$ spin-orbitals, shells
$0$--$5$) for $N=2,6,12$ electrons, with ancillas to spare. The limiter is
\textbf{circuit depth}, and these runs will \emph{validate} the pipeline rather
than beat classical methods.
\end{block}
\vspace{2pt}
\begin{itemize}
  \item Mapping here is \textbf{1 qubit per spin-orbital} (not 2/site as in Hubbard):
        $K$ spin-orbitals $\to K$ qubits.
  \item \textbf{$\sim$12 qubits} already gives a non-trivial, exactly-checkable
        correlation benchmark (the notebook's $K=12$ case).
  \item \textbf{$\sim$42 qubits} reaches the full shells-$0$--$5$ space --- a
        $\sim\!10^{10}$-determinant FCI, right at the classical exact-FCI frontier.
  \item No quantum advantage: coupled cluster / DMRG / selected-CI own this regime.
        Value $=$ \textbf{clean pipeline validation} and high-overlap state prep for Rodeo.
\end{itemize}
\end{frame}

% ---------- mapping ----------
\begin{frame}{The encoding: $K$ spin-orbitals $\to K$ qubits}
\begin{columns}[T]
\begin{column}{0.54\textwidth}
Each harmonic-oscillator shell $N_s=2n+|m|$ holds $2(N_s{+}1)$ spin-orbitals.
Under Jordan--Wigner, \textbf{one spin-orbital $=$ one qubit}.
\vspace{4pt}
\begin{center}\small
\renewcommand{\arraystretch}{1.15}
\begin{tabular}{@{}ccccc@{}}
\toprule
shells & spin-orb.\ $K$ & \textbf{qubits} & magic $N$ \\
\midrule
$0$       & 2  & 2  & 2 \\
$0$--$1$  & 6  & 6  & 6 \\
$0$--$2$  & 12 & 12 & 12 \\
$0$--$3$  & 20 & 20 & 20 \\
$0$--$4$  & 30 & 30 & 30 \\
$0$--$5$  & \textbf{42} & \textbf{42} & 42 \\
$0$--$6$  & 56 & \alert{56} & 56 \\
\bottomrule
\end{tabular}
\end{center}
\end{column}
\begin{column}{0.46\textwidth}
\begin{block}{Headroom on Magne}
50 logical qubits $\Rightarrow$ active space up to \textbf{$K\approx47$--$49$}
spin-orbitals (after $1$--few Rodeo ancillas).
\begin{itemize}\small
  \item Shells $0$--$5$ ($K=42$): \textbf{fits comfortably}.
  \item Shells $0$--$6$ ($K=56$): \textbf{does not fit}.
\end{itemize}
So the natural maximal dot active space on Magne is \textbf{$K=42$}.
\end{block}
\end{column}
\end{columns}
\end{frame}

% ---------- qubit budget table ----------
\begin{frame}{Qubit budget vs.\ classical cost}
\begin{center}\small
\renewcommand{\arraystretch}{1.18}
\begin{tabular}{@{}rrlrl@{}}
\toprule
$N$ & \textbf{qubits $K$} & shells & FCI dim $\binom{K}{N}$ & classical status \\
\midrule
2  & 12 & $0$--$2$ & 66            & trivial \\
2  & 20 & $0$--$3$ & 190           & trivial \\
2  & 42 & $0$--$5$ & 861           & trivial \\
6  & 20 & $0$--$3$ & 38{,}760      & trivial \\
6  & 30 & $0$--$4$ & $5.9\times10^{5}$ & trivial \\
6  & 42 & $0$--$5$ & $5.2\times10^{6}$ & easy exact FCI \\
12 & 30 & $0$--$4$ & $8.6\times10^{7}$ & easy exact FCI \\
12 & 42 & $0$--$5$ & $1.1\times10^{10}$ & \textbf{exact-FCI frontier} \\
20 & 42 & $0$--$5$ & $5.1\times10^{11}$ & beyond exact FCI (CC/DMRG ok) \\
\bottomrule
\end{tabular}
\end{center}
\vspace{2pt}
{\small\color{qngray} Symmetry sectors ($M$, $S_z$) shrink these $\sim 5$--$20\times$.
Note: at the magic numbers $N$ sits far from half-filling, so FCI dimensions stay
modest --- the dot is \emph{not} a worst-case correlation problem.}
\end{frame}

% ---------- classical frontier ----------
\begin{frame}{Where the classical frontier is for quantum dots}
\begin{itemize}
  \item Exact FCI is feasible to $\binom{K}{N}\!\sim\!10^{10}$ determinants ---
        i.e.\ up to the $N{=}12$, $K{=}42$ case \emph{on a workstation/cluster}.
  \item Beyond that, \textbf{coupled cluster (CCSD(T)), DMRG, selected-CI and AFQMC}
        are accurate \emph{and cheap} for these systems.
  \item Quantum dots in this notebook's regime are \textbf{smooth and
        moderately correlated}: HF is a good reference, correlation converges by
        $K\!\approx\!12$, UCCSD recovers it in a $\sim$7-parameter landscape.
\end{itemize}
\begin{alertblock}{Consequence}
There is \textbf{no quantum-dot regime reachable with 50 logical qubits that beats
classical methods.} This is expected --- and fine. The scientific role of these runs
is \textbf{algorithm validation against an exact reference}, and demonstrating
high-overlap initial states for Rodeo.
\end{alertblock}
\end{frame}

% ---------- how many qubits for serious ----------
\begin{frame}{So how many logical qubits for a ``serious'' benchmark?}
\renewcommand{\arraystretch}{1.3}
\begin{center}\small
\begin{tabular}{@{}p{3.2cm}p{2.2cm}p{6.4cm}@{}}
\toprule
\textbf{Benchmark tier} & \textbf{qubits} & \textbf{What it demonstrates} \\
\midrule
Sanity / shakedown & 2--6 & JW Hamiltonian, $Z$-basis readout, HF state prep. \\
\textbf{Meaningful} & \textbf{12} & First real correlation (notebook's $K{=}12$,
   $\sim$7 UCCSD amplitudes); exact FCI check; full VQE$\to$Rodeo loop. \\
Substantial & 20--30 & $N{=}2,6$ in larger bases; non-trivial UCCSD depth;
   approaches basis-set convergence. \\
Maximal natural & 42 & Shells $0$--$5$; $N{=}12$ is a $\sim\!10^{10}$-det FCI ---
   a genuinely meaty benchmark, \emph{if} depth allows. \\
\bottomrule
\end{tabular}
\end{center}
\vspace{2pt}
\begin{block}{Verdict on the qubit axis}
\textbf{12} qubits for a credible first benchmark, \textbf{20--30} for a substantial one,
\textbf{42} for the largest natural dot space --- all inside 50.
\textbf{Qubit count is not the limiter.}
\end{block}
\end{frame}

% ---------- depth is the limiter ----------
\begin{frame}{The binding constraint is depth (again)}
Raw Trotterised UCCSD two-qubit-gate counts under Jordan--Wigner (1 Trotter step):
\begin{center}\small
\renewcommand{\arraystretch}{1.15}
\begin{tabular}{@{}rrrrr@{}}
\toprule
$N$ & $K$ & singles & doubles & $\sim$2-qubit gates \\
\midrule
2  & 12 & 20  & 45     & $\sim\!10^{4}$ \\
6  & 20 & 84  & 1{,}365 & $\sim\!4\times10^{5}$ \\
12 & 30 & 216 & 10{,}098 & $\sim\!5\times10^{6}$ \\
12 & 42 & 360 & 28{,}710 & $\sim\!2\times10^{7}$ \\
\bottomrule
\end{tabular}
\end{center}
\begin{itemize}\small
  \item JW adds non-local $Z$-strings of length up to $K$ per term --- neutral-atom
        reconfigurable connectivity helps, but depth still grows fast.
  \item Rodeo adds $M$ rounds of controlled-$e^{-iH t_j}$; $H$ has
        $\sim\!10^{2}$ ($K{=}12$) to $\sim\!10^{5}$ ($K{=}42$) Pauli terms per Trotter step.
  \item \textbf{Depth-feasible on an early Level-2 machine: $K\!=\!12$, perhaps $K\!=\!20$
        with ADAPT-VQE.} $K\!=\!30$--$42$ are qubit-feasible but \emph{depth-prohibitive} at first.
\end{itemize}
\end{frame}

% ---------- recommendations ----------
\begin{frame}{Recommendations}
\begin{columns}[T]
\begin{column}{0.5\textwidth}
\textbf{Make depth, not width, the design driver}
\begin{itemize}\small
  \item Use \textbf{ADAPT-VQE} / symmetry-adapted pools: the raw $28{,}710$ doubles at
        $K{=}42$ collapse to $O(10$--$100)$ selected operators.
  \item Exploit the \textbf{division of labour}: VQE need only reach overlap
        $p\!\approx\!0.9$--$0.95$ --- \emph{Rodeo} supplies the precision. A \emph{shallow}
        UCCSD that gets high overlap is enough; do not over-optimise VQE depth.
\end{itemize}
\end{column}
\begin{column}{0.5\textwidth}
\textbf{Choose friendly physics \& honest claims}
\begin{itemize}\small
  \item Work at \textbf{moderate $\hbar\omega$} (weak/intermediate coupling) where HF is a
        good reference; avoid the small-$\hbar\omega$ Wigner-molecule regime that needs deeper circuits.
  \item In the application, replace ``20--40 qubits'' with explicit $(N,K)$:
        e.g.\ \emph{``$N{=}2$--$12$ dots in active spaces up to $K{=}42$ spin-orbitals
        ($\le 42$ qubits).''}
  \item Always pair a run with its \textbf{exact-FCI reference} --- the dot's great virtue.
\end{itemize}
\end{column}
\end{columns}
\end{frame}

% ---------- benchmark ladder ----------
\begin{frame}{A concrete benchmark ladder for Magne}
\begin{enumerate}
  \item \textbf{$K=2$ ($N=2$)} --- hardware sanity: HF state, diagonal $H$, Rodeo on a 1-state space.
  \item \textbf{$K=6$ ($N=2$)} --- first off-diagonal scattering; validate VQE energy vs.\ FCI ($\dim 15$).
  \item \textbf{$K=12$ ($N=2$) --- the flagship first benchmark.} $\sim$7 UCCSD amplitudes,
        FCI $\dim 66$, full VQE$\to$Rodeo with exact reference. \emph{$\sim$12 logical qubits.}
  \item \textbf{$K=20$ ($N=2$ or $6$)} --- ADAPT-VQE; measure rodeo convergence vs.\ rounds $M$.
        \emph{$\sim$20 qubits.}
  \item \textbf{Stretch: $K=30$--$42$ ($N=6,12$)} --- only as the logical depth budget allows;
        $N{=}12,K{=}42$ is the $\sim\!10^{10}$-det showcase. \emph{30--42 qubits.}
\end{enumerate}
\vspace{2pt}
{\small\color{qngray} Every rung has an exact classical answer, so each one cleanly
separates \emph{algorithm error} from \emph{hardware error} --- exactly what a Level-2
demonstration needs.}
\end{frame}

% ---------- closing ----------
\begin{frame}{Closing summary}
\begin{block}{Feasibility of quantum dots on Magne (50 logical qubits)}
\textbf{Qubit-feasible across the board:} 12 qubits for a meaningful benchmark,
20--30 for a substantial one, 42 for the full shells-$0$--$5$ space --- all within budget.
\textbf{Depth-limited:} only the smaller cases ($K\!\le\!20$, with ADAPT-VQE) will run to
convergence on an early Level-2 machine. \textbf{No classical advantage} --- and none expected.
\end{block}
\begin{itemize}
  \item The quantum dot is the \textbf{ideal validation system}: exact FCI reference at every size,
        smooth correlation, and a clean test of VQE--UCCSD $\to$ Rodeo end-to-end.
  \item Spend the qubits you have on \textbf{depth-reducing structure} (ADAPT, symmetry,
        shallow high-overlap VQE), not on the largest possible basis.
  \item Use it to \textbf{characterise Magne and de-risk the pipeline} before the
        harder Hubbard / nuclear targets.
\end{itemize}
\end{frame}

\end{document}
